\chapter{Captura de requisitos y construcción del benchmark experimental}
\label{cap:capitulo3}

Este capítulo describe el proceso seguido para definir los requisitos funcionales del sistema y diseñar el conjunto de pruebas utilizado posteriormente para evaluar las capacidades de los modelos de lenguaje analizados. Dado que el objetivo principal del proyecto consiste en facilitar la consulta de colecciones musicales por parte de investigadoras del ámbito musicológico, la identificación de necesidades reales de uso constituyó una etapa fundamental del desarrollo.

Para ello, se llevó a cabo un proceso de captura de requisitos con musicólogas mediante entrevistas personales y formularios estructurados. Estas actividades permitieron identificar los tipos de consultas más frecuentes, las dificultades encontradas al trabajar con grandes corpus musicales y las necesidades analíticas que una herramienta basada en modelos de lenguaje debería satisfacer.

A partir de la información recopilada se construyó un conjunto de preguntas de evaluación o \textit{benchmark}, diseñado para reproducir escenarios reales de investigación musicológica. Este \textit{benchmark} constituye la base experimental empleada en los capítulos posteriores para comparar distintos enfoques tecnológicos y analizar sus fortalezas y limitaciones en tareas de consulta, recuperación y análisis de información musical.

\section{Captura de requisitos con musicólogas a través de entrevistas personales}
% Para recopilar información sobre las necesidades, expectativas y posibles usos de la herramienta desarrollada, se realizaron entrevistas presenciales con musicólogas de la Facultad de Geografía e Historia de la Universidad Complutense de Madrid. Estas entrevistas permitieron obtener información cualitativa más detallada sobre los criterios de búsqueda, análisis y contextualización considerados relevantes en el ámbito de la musicología y el estudio de la música tradicional y folclórica.

Antes de abordar el diseño e implementación de la herramienta, se consideró necesario identificar las necesidades reales de las potenciales usuarias finales. Aunque las posibilidades técnicas ofrecidas por los modelos de lenguaje y las tecnologías de recuperación de información son amplias, la utilidad práctica de un sistema de estas características depende en gran medida de su capacidad para responder a problemas concretos presentes en los flujos de trabajo de investigación. Por este motivo, se llevó a cabo una fase inicial de captura de requisitos orientada a comprender cómo las investigadoras formulan sus preguntas, qué tipo de información buscan habitualmente y cuáles son las principales limitaciones que encuentran al trabajar con colecciones musicales de gran tamaño.

Dado que el proyecto se encuentra orientado al análisis y exploración de repertorios musicales tradicionales, las musicólogas especializadas en el estudio de la música folclórica y del patrimonio musical constituyen el perfil de usuario más representativo para evaluar la utilidad potencial de la herramienta. Estas investigadoras poseen experiencia directa en tareas de catalogación, consulta de cancioneros, análisis de partituras y contextualización histórico-cultural de repertorios musicales, por lo que su conocimiento resulta especialmente valioso para identificar funcionalidades relevantes desde una perspectiva aplicada.

Con este objetivo, se realizaron entrevistas presenciales con musicólogas de la Facultad de Geografía e Historia de la Universidad Complutense de Madrid. La elección de entrevistas permitió obtener información cualitativa de mayor profundidad que la proporcionada por instrumentos cerrados de recogida de datos, favoreciendo la aparición de comentarios espontáneos, necesidades no previstas inicialmente y ejemplos concretos extraídos de la práctica investigadora cotidiana. Asimismo, este formato facilitó la discusión de posibles escenarios de uso y permitió comprender mejor los procesos de trabajo empleados actualmente para consultar, analizar y comparar información musical.

Las entrevistas se centraron en identificar los criterios de búsqueda más frecuentes, los tipos de consultas consideradas de mayor interés musicológico, las dificultades encontradas al trabajar con grandes corpus digitales y el grado de utilidad percibido de una interfaz basada en lenguaje natural. La información recopilada durante esta fase no solo contribuyó a orientar el diseño funcional de la herramienta desarrollada, sino que también sirvió como base para la construcción del conjunto de preguntas de evaluación (\textit{benchmark}) utilizado posteriormente para analizar las capacidades de los distintos modelos de lenguaje estudiados.

Entre los principales intereses identificados destacó la necesidad de realizar búsquedas musicales avanzadas a partir de elementos concretos de las obras, como los íncipits musicales (los primeros compases instrumentales o vocales), estructuras melódicas, letras, instrumentación o características interpretativas. Asimismo, las investigadoras señalaron la importancia de disponer de información musicológica detallada, incluyendo aspectos como autoría, fecha, duración, tonalidad, tempo, rango vocal o presencia de grabaciones sonoras.

Otro de los aspectos recurrentes fue el interés por analizar la evolución y transmisión de las obras a lo largo del tiempo, especialmente en el contexto de la música tradicional y folclórica. En este sentido, se destacó la relevancia de estudiar las variaciones entre versiones, los procesos de creación y transmisión oral, así como las relaciones existentes entre músicas de distintas regiones, comunidades y contextos culturales.

Las entrevistas también pusieron de manifiesto la importancia de incorporar funcionalidades orientadas a la comparación y análisis relacional de obras musicales, tales como la identificación de patrones compartidos, la frecuencia de aparición de determinados íncipits o leitmotivs musicales dentro de un corpus o la comparación simultánea de canciones dentro de una misma interfaz.

Asimismo, las participantes señalaron la necesidad de integrar información procedente de diferentes catálogos y bases de datos especializadas (como BNE, ISNI, Wikidata o SGAE), incluyendo variantes de nombres y pseudónimos de autores e intérpretes con el fin de facilitar procesos de recuperación y contextualización de la información musical.

Finalmente, las entrevistas evidenciaron el valor de conservar metadatos relacionados con la recopilación y procedencia de las obras, tales como el lugar de recuperación, el contexto de documentación o los datos asociados al proceso de recogida de la información, aspectos especialmente relevantes en investigaciones vinculadas al patrimonio musical y la tradición oral.

\section{Captura de requisitos con musicólogas a través de formularios}
Además de las entrevistas realizadas de manera presencial, se diseñó un formulario estructurado dirigido a potenciales usuarias pertenecientes al ámbito académico e investigador. El objetivo principal de este instrumento fue identificar los perfiles de los participantes, conocer los tipos de consultas que realizarían con la herramienta y comprender el uso previsto de los resultados obtenidos.
El cuestionario fue creado y difundido mediante Google Forms, lo que permitió una distribución sencilla y accesible, así como la recopilación organizada de las respuestas.

El cuestionario incluyó, en primer lugar, una pregunta destinada a caracterizar el perfil académico de los participantes, ofreciendo las categorías de investigador/a, estudiante, profesor/a u otras situaciones profesionales relacionadas. Asimismo, se solicitó información sobre el área de investigación de cada encuestada con el fin de contextualizar sus respuestas y analizar posibles diferencias entre disciplinas.

A continuación se incorporaron dos preguntas abiertas orientadas a identificar las consultas que las usuarias podrían considerar relevantes a realizar mediante la herramienta y a documentar el uso que las académicas darían a los datos proporcionados por el sistema tras realizar dicha consulta. 
La información recopilada permitió, por una parte, recoger ejemplos concretos de necesidades de búsqueda, análisis o recuperación de información, así como otros posibles escenarios de uso no previstos inicialmente por los investigadores y, por otra parte, explorar el potencial impacto de la herramienta en actividades de investigación, docencia, análisis, documentación o generación de nuevo conocimiento.

Las respuestas obtenidas mediante el formulario permitieron identificar distintos intereses y necesidades relacionadas con la investigación musicológica y el desarrollo de herramientas de recuperación y análisis de información musical. Las participantes pertenecían principalmente al ámbito investigador y docente, con especialización en áreas como la música latinoamericana de los siglos XIX y XX, la producción musical y el teatro lírico, así como la musicología y los estudios sobre flamenco y folclore.

Uno de los aspectos más destacados fue el interés por disponer de funcionalidades de búsqueda avanzadas que permitan localizar distintas versiones de una misma obra, clasificar repertorios por estilos o estructuras musicales y acceder a diferentes interpretaciones y grabaciones relacionadas con un mismo material musical. Asimismo, varias respuestas señalaron la importancia de incorporar conexiones con otros recursos externos, como archivos, bibliotecas o bases de datos especializadas, con el fin de contextualizar las obras y establecer relaciones entre distintos repertorios musicales y literarios.

En el ámbito de la música tradicional y el folclore, las respuestas pusieron de manifiesto la relevancia de incluir criterios de búsqueda relacionados con aspectos temáticos, geográficos y socioculturales. Entre ellos destacan la posibilidad de identificar referencias a lugares concretos, analizar temáticas recurrentes en las letras o estudiar variaciones lingüísticas y culturales en función de la región, el idioma o la época histórica.

Asimismo, las musicólogas señalaron el interés de utilizar los datos recuperados para la construcción de bases de datos propias, el desarrollo de proyectos de investigación, la elaboración de materiales docentes y la creación de repertorios especializados. También se destacó el potencial de la herramienta para generar visualizaciones y mapas interactivos que permitan representar la distribución geográfica y la evolución de determinadas tradiciones musicales.

Otro aspecto recurrente fue la importancia atribuida a la documentación del proceso de recopilación y transmisión de la música tradicional. Se consideró relevante incorporar información sobre la labor desarrollada por folcloristas y etnomusicólogos, incluyendo datos sobre las regiones estudiadas, los repertorios recopilados y las metodologías de recogida empleadas.

En conjunto, los resultados del formulario evidencian la necesidad de una herramienta flexible y adaptada, capaz de integrar información procedente de múltiples fuentes y de ofrecer funcionalidades avanzadas de búsqueda, análisis y contextualización musical. Las respuestas obtenidas refuerzan además la importancia de abordar la música no únicamente como un objeto textual o sonoro aislado, sino también como un fenómeno cultural, histórico y social.

\section{Formulación de preguntas para comprobar las capacidades de los modelos de lenguaje}
% Como parte del diseño experimental inicial, se elaboró un conjunto de preguntas teórico-musicales orientadas a evaluar las capacidades potenciales de los modelos de lenguaje pensados para la herramienta en distintos ámbitos de la musicología, el análisis musical y la recuperación avanzada de información musical. El objetivo de este proceso fue definir un repertorio amplio y representativo de consultas que permitiera cubrir tanto aspectos estrictamente musicales —como ritmo, tonalidad, métrica o melodía— como dimensiones documentales, textuales, históricas y técnicas relacionadas con los formatos de codificación musical utilizados en el corpus.

% Las preguntas fueron formuladas con distintos niveles de complejidad y especificidad, combinando criterios musicales, metadatos descriptivos, análisis lírico y aspectos estructurales derivados de formatos como MEI, MusicXML o MSCZ. De este modo, se buscó evaluar la viabilidad de consultas complejas capaces de relacionar múltiples dimensiones analíticas dentro de un mismo entorno de recuperación de información.

% Para facilitar su organización y análisis, las preguntas fueron agrupadas en diferentes bloques temáticos en función del ámbito musicológico o técnico abordado.

El diseño de las preguntas utilizadas para la evaluación de los modelos de lenguaje se fundamentó directamente en los requisitos identificados durante las entrevistas y formularios realizados con las potenciales usuarias de la herramienta. De este modo, se buscó garantizar que el \textit{benchmark} reprodujera escenarios reales de investigación musicológica y no únicamente ejemplos teóricos diseñados desde una perspectiva tecnológica.

A partir del análisis de las respuestas obtenidas se identificaron cinco grandes categorías funcionales: recuperación de información musical básica, análisis musical avanzado y comparación estructural, relación entre texto y música, contextualización musicológica y etnomusicológica y análisis de estructuras técnicas de codificación musical. Cada una de estas categorías fue posteriormente traducida a un conjunto específico de preguntas destinadas a evaluar la capacidad de los modelos para resolver tareas representativas de dicho requisito.

Este enfoque permitió establecer una trazabilidad explícita entre las necesidades expresadas por las usuarias y los experimentos realizados posteriormente, reforzando la validez metodológica del proceso de evaluación y facilitando la interpretación de los resultados obtenidos.

\begin{itemize}
\item \textbf{Ritmo, métrica y tempo.} Preguntas centradas en la acentuación, figuras rítmicas, síncopas, polirritmias, velocidad (BPM) y la relación matemática/temporal de las notas dentro del compás. Preguntas incluidas:
    \begin{itemize}
    \item \textit{Identifica piezas en 2/4 que contengan síncopas ligadas entre el segundo pulso de un compás y el primero del siguiente de forma sistemática}
    \item \textit{Compara la densidad de notas por compás entre una pieza de 4/4 y una de 2/4 del dataset; ¿cuál tiene mayor promedio de eventos musicales?}
    \item \textit{Identifica si existe algún uso de polirritmia en los archivos MEI}
    \item \textit{Identifica si existe algún patrón de silencios (<rest />) que se repita sistemáticamente cada cuatro compases en las piezas de métrica ternaria}
    \item \textit{Identifica piezas con un compás de 4/4}
    \item \textit{Identifica piezas con hemiolias horizontales}
    \item \textit{Identifica piezas con hemiolias verticales}
    \item \textit{Identifica piezas con un bpm de 120}
    \item \textit{Identifica piezas con síncopas}
    \item \textit{Identifica piezas con polirritmia}
    \end{itemize}

\item \textbf{Análisis melódico y armónico.} Consultas orientadas a los intervalos musicales, tesituras, movimientos melódicos (ascendentes/descendentes), alteraciones accidentales (sostenidos, bemoles y becuadros), tonalidades y grados de resolución (tónica/dominante). Preguntas incluidas:
    \begin{itemize}
    \item \textit{Localiza todas las notas que tengan bemoles accidentales y dime en qué compases aparecen en los archivos MEI encontrados}
    \item \textit{¿Cuál es el intervalo melódico más frecuente entre la última nota de un compás y la primera del siguiente en las piezas de modo mixolidio?}
    \item \textit{Identifica piezas donde la nota final sea la dominante (quinto grado) en lugar de la tónica, sugiriendo una estructura abierta}
    \item \textit{Identifica piezas con la tonalidad G major}
    \item \textit{Identifica piezas que estén en modo dórico}
    \item \textit{Identifica piezas cuya nota más grave sea A3}
    \item \textit{Identifica piezas cuya nota más aguda sea G5}
    \item \textit{Identifica obras cuya tendencia melódica sea predominantemente descendente}
    \item \textit{Identifica piezas con 4 alteraciones accidentales}
    \item \textit{Identifica piezas cuyo intervalo frontera más frecuente sea +5}
    \end{itemize}

\item \textbf{Relación métrica, lírica y texto.} Preguntas híbridas que analizan cómo interactúa el texto (la lírica, el número de sílabas, palabras clave) con los aspectos musicales del tempo, la altura de las notas o la duración rítmica. Preguntas incluidas:
    \begin{itemize}
    \item \textit{Identifica canciones donde el autor sea una mujer, pero el contenido lírico esté escrito en voz masculina}
    \item \textit{Filtra piezas que usen la palabra ``estrella'' en la lírica y determina si su melodía es predominantemente ascendente o descendente}
    \item \textit{Compara el uso de las sílabas ``duer'' y ``me'' en los distintos archivos del dataset: ¿se asigna siempre a la misma duración de nota?}
    \item \textit{Determina si las canciones con tempo superior a 120 BPM tienden a tener letras con palabras de menor longitud silábica}
    \item \textit{Identifica piezas cuyo contenido lírico esté escrito en voz masculina}
    \item \textit{Identifica piezas que usen la palabra ``bondad''}
    \item \textit{Identifica piezas que mencionen el nombre propio ``Alfonsito''}
    \end{itemize}

\item \textbf{Contexto sociocultural, temático y geográfico.} Preguntas complejas que cruzan la etnomusicología o el análisis musical comparativo con datos del entorno (temas como ``Nanas'' o ``Trabajo'' o fechas de recolección históricas). Preguntas incluidas:
    \begin{itemize}
    \item \textit{Busca ejemplos de anacrusa en piezas que estén etiquetadas como ``Nanas''}
    \item \textit{Encuentra piezas recolectadas antes de 1950 que traten el tema de ``la muerte'' y que utilicen exclusivamente un rango melódico de sexta menor}
    \item \textit{¿Qué diferencias de tesitura (nota más aguda vs. nota más grave) existen entre las canciones de nana?}
    \item \textit{Busca canciones que traten temas de trabajo que compartan la misma estructura de intervalos que las canciones de tema infantil}
    \item \textit{Identifica piezas que se llamen El Carro}
    \item \textit{Identifica piezas cuyo autor sea Schumann}
    \item \textit{Identifica piezas cuyo autor sea de género masculino}
    \item \textit{Identifica piezas que traten la muerte como tema}
    \item \textit{Identifica piezas de tonalidad A minor que tengan hemiolias y que no superen los 90 bpm}
    \item \textit{Identifica piezas de tonalidad A minor que tengan hemiolias}
    \item \textit{Lista canciones que compartan la misma armadura de un bemol (F major/D minor) pero que traten temas distintos}
    \end{itemize}

\item \textbf{Estructura analítica de datos y formatos técnicos.} Este bloque engloba las preguntas que se centran puramente en el código o metadatos de los archivos musicales (MEI, MusicXML, XML, MSCZ), discrepancies de sintaxis musical y atributos técnicos de renderizado (como plicas, volúmenes MIDI o transcripciones). Preguntas incluidas:
    \begin{itemize}
    \item \textit{Encuentra canciones donde el valor de dur.ppq cambie a mitad de la sección, indicando un cambio de resolución rítmica}
    \item \textit{Localiza obras que tengan compases donde la suma de las duraciones de las etiquetas <note> no coincida con el meterSig declarado en el staffDef}
    \item \textit{¿Qué piezas utilizan valores irregulares (por ejemplo, tresillos) no declarados explícitamente en un <time-modification> del MusicXML?}
    \item \textit{Identifica el uso de la etiqueta <sb /> (system break) y determina si coincide siempre con el final de una frase musical lógica}
    \item \textit{Identifica piezas con un volumen midi de 79}
    \item \textit{Identifica piezas cuyo software de codificación haya sido Sibelius 7.1.3}
    \item \textit{Identifica piezas cuya fecha de codificación sea el 2015-07-01}
    \item \textit{Identifica piezas con compases vacíos}
    \item \textit{Identifica piezas con notas de duración 0}
    \item \textit{Identifica piezas con plicas anómalas}
    \end{itemize}
\end{itemize}